\section{Stability of the maintained optical state}\label{sec:stability}

\subsection{A stationary force balance is only the starting point}
Let $(\bm y_\star,\bm a_\star)$ be a candidate maintained state and command.
For an autonomous slow-state or acoustic-envelope description, define
$\bm\chi=\bm y-\bm y_\star$ in local coordinates satisfying the conservation
constraints. Denote command perturbation by $\delta\bm a=\bm a-\bm a_\star$,
and the dynamic state and command derivative operators by $\mathsf A$ and
$\mathsf B$, respectively. Linearization gives
\begin{equation}
 \dot{\bm\chi}=\mathsf A\bm\chi+\mathsf B\delta\bm a.
 \label{eq:linear-state}
\end{equation}
The operators $\mathsf A,\mathsf B$ are derivatives of the dynamic equations,
not derivatives of a surface-fitting routine. The state includes the necessary
fluid, thermal, actuator and acoustic memory variables. A steady phasor field
represents a periodic physical wave; if an envelope reduction is unjustified,
stability must instead be studied about that periodic state.

For a finite-dimensional autonomous reduction, negative real parts of all
eigenvalues imply local linear asymptotic stability. For the continuum problem,
point spectrum alone is insufficient: suitable evolution-operator bounds are
needed. Define $\mathcal T(t)$ as the linear perturbation propagation operator,
$M_s\geq1$ as a dimensionless amplification bound, and $\alpha_s>0$ as a decay
rate. A stronger useful condition is
\begin{equation}
 \norm{\mathcal T(t)}\leq M_s e^{-\alpha_s t},\qquad t\geq0.
 \label{eq:semigroup-bound}
\end{equation}
The norm uses physically stated state scales. A large $M_s$ permits substantial
transient growth despite eventual decay. Optical acceptability additionally
depends on the observation $\mathsf C_{\rm opt}\bm\chi$, where
$\mathsf C_{\rm opt}$ is the linearized optical measurement operator, and on
the amplitude of admitted disturbances. Numerical convergence, shape accuracy
and dynamical stability are three different claims.

\subsection{Capillary stiffness and acoustic feedback}
In the quiescent, constant-$\sigma$ reduction, define the shape operator
$\bm B_\Gamma=\Gs\nn$ and its squared norm $|\bm B_\Gamma|^2$, equal to the sum
of squared principal curvatures. Let $\Delta_\Gamma=\Gs\cdot\Gs$ denote the
surface Laplacian. Under normal displacement $\eta$,
\begin{equation}
 \delta\kappa=-(\Delta_\Gamma+|\bm B_\Gamma|^2)\eta,
 \qquad \delta\Phi=(\nn\cdot\nabla\Phi)\eta.
 \label{eq:curvature-variation}
\end{equation}
At fixed actual commanded variable, denote the complete shape derivative of
acoustic normal load by $D_\Gamma\Pi[\eta]$. The restoring-pressure operator,
before projection onto volume-preserving variations, is
\begin{equation}
 \boxed{\mathcal K_{\rm eff}\eta
 =-\sigma(\Delta_\Gamma+|\bm B_\Gamma|^2)\eta
    +\Delta\rho(\nn\cdot\nabla\Phi)\eta
    -D_\Gamma\Pi[\eta].}
 \label{eq:effective-stiffness}
\end{equation}
The pressure-multiplier variation is removed on that constrained space. Boundary
terms and contact-line perturbations must be included for the chosen wetting
model. Thermal variations add derivatives of $\sigma$, density and other
properties. This formula shows the relevant question: when the interface moves,
does the resulting acoustic load counteract or reinforce the displacement?
Matching the load at one geometry supplies no answer to that derivative.

If fluid inertia and acoustic memory are negligible, let $\mathcal M_\Gamma$
be the viscous mobility mapping outward traction to interface normal velocity,
with its constraints and boundary conditions. Then the creeping-flow limit is
\begin{equation}
 \dot\eta=-\mathcal M_\Gamma\mathcal K_{\rm eff}\eta.
 \label{eq:stokes-stability}
\end{equation}
For passive Stokes flow the constrained mobility is positive in the reciprocal
work pairing. A coercive self-adjoint restoring operator can then give an
energy-based stability criterion. Driven, dissipative acoustic feedback need
not be self-adjoint or derive from a scalar shape potential; positivity of a
guessed ``capillary plus acoustic energy'' is not a general stability proof.

With memory retained, let $s\in\C$ be a Laplace frequency,
$\mathcal Z_h(s)$ the hydrodynamic restoring/dynamic impedance, and
$\mathcal K_a(s)$ the linearized acoustic load response to surface displacement.
The homogeneous perturbation condition is
\begin{equation}
 [\mathcal Z_h(s)-\mathcal K_a(s)]\eta=0.
 \label{eq:memory-stability}
\end{equation}
Both operators include their relevant boundary and state coupling. In a
particular modal approximation, hydrodynamics may be represented by
$\mathcal Z_h(s)=s^2\mathsf M+s\mathsf D+\mathsf K_h$, where
$\mathsf M,\mathsf D,\mathsf K_h$ are the modal mass, damping and passive
stiffness matrices. Frequency-independent damping is an approximation, not a
universal viscous free-surface law; capillary-wave dispersion depends on the
viscous regime \cite{denner2016}. Eliminating acoustic state instantaneously
replaces $\mathcal K_a(s)$ by its zero-frequency value and can miss an unstable
oscillatory mode.

For a periodically maintained state of period $T_p$, let $\mathcal M_p$ denote
the linear evolution over one period. In a finite-dimensional periodic
reduction, its Floquet multipliers must lie strictly inside the unit circle
after genuine symmetry modes are treated appropriately. A continuum statement
also requires operator control. Pulse trains, modulation and beat forcing may
therefore need a periodic analysis even when their mean load is correct.

\subsection{Feedback must have the necessary authority and observations}
Let the measured perturbation be
$\bm z_{\rm obs}=\mathsf C\bm\chi+\bm\nu$, where $\mathsf C$ is the sensor
operator and $\bm\nu$ measurement noise. The existence of a desired equilibrium
does not imply that its unstable modes can be controlled or observed. Local
stabilization requires, at minimum, authority over the unstable modes and
sufficient information to estimate those that need feedback. Delays, saturation
and finite sensing bandwidth belong to the control model.

For illustration only, consider real, scaled finite-dimensional coordinates
and full-state feedback $\delta\bm a=-\mathsf K_c\bm\chi$. Here
$\mathsf K_c$ is a feedback gain, distinct from the traction matrices
$\bm K_i$. Let $\mathsf P_L\succ0$ be a symmetric Lyapunov matrix. A
sufficient linear decay condition is
\begin{equation}
 (\mathsf A-\mathsf B\mathsf K_c)^T\mathsf P_L
 +\mathsf P_L(\mathsf A-\mathsf B\mathsf K_c)
 \preceq-2\alpha_s\mathsf P_L.
 \label{eq:lyapunov}
\end{equation}
With $\mathsf X=\mathsf P_L^{-1}$ and $\mathsf Y=\mathsf K_c\mathsf X$, this
becomes the linear matrix inequality
\begin{equation}
 \mathsf A\mathsf X+\mathsf X\mathsf A^T
 -\mathsf B\mathsf Y-\mathsf Y^T\mathsf B^T
       +2\alpha_s\mathsf X\preceq0,\qquad\mathsf X\succ0.
 \label{eq:feedback-lmi}
\end{equation}
This established construction \cite{boyd1994} is convex for fixed
$\mathsf A,\mathsf B,\alpha_s$ and the stated full-state model. Joint design
of geometry, nonlinear acoustics and output feedback does not inherit that
convexity. A controller satisfying it still needs a compatible estimator,
actuator headroom and a nonlinear region of attraction containing the intended
formation endpoint and disturbances.
